% Root file used ashow macros
\documentclass{beamer}
\usepackage{tikz}
\usepackage{multicol}

% ================================================================
% Answer-overlay helpers
% ================================================================
\newcommand{\ablank}[1]{%
  \alt<2>{\textcolor{red}{#1}}{\underline{\phantom{#1}}}%
}
\newcommand{\ashow}[1]{\uncover<2->{\textcolor{red}{\small #1}}}
\newcommand{\ashowq}[1]{\alt<2->{\textcolor{red}{\small #1}}{?}}

% Answers drawn straight on to a TikZ diagram, revealed on overlay 2 like \ashow.
% Space is reserved on overlay 1, so the diagram does not move between slides.
\usetikzlibrary{overlay-beamer-styles}
\tikzset{
  ans/.style     = {red, thick, visible on=<2->},
  ansfill/.style = {red!15, visible on=<2->},
  anslab/.style  = {red, font=\tiny, visible on=<2->},
}

% ---- Minimal styling ----
\setbeamertemplate{navigation symbols}{}
\setbeamertemplate{footline}{}
\setbeamercolor{background canvas}{bg=white}
\setbeamercolor{frametitle}{fg=black}
\setbeamerfont{frametitle}{size=\Large, series=\bfseries}
\usefonttheme{professionalfonts}

% ---- Title information for \frame{\titlepage} ----
\title{Sequences Starters}
\author{}
\date{}

% ---- Shared tikz ----
\tikzset{every picture/.style={baseline=(current bounding box.south)}}

% 1) matchstick squares (n squares)
\newcommand{\matchsquares}[2][black]{%
  \begin{tikzpicture}[scale=0.65]
    \pgfmathsetmacro\n{#2}
    \color{#1}
    \foreach \x in {0,...,\n} {
      \draw[very thick, cap=round] (\x,0) -- (\x,1);
    }
    \foreach \x in {0,...,\the\numexpr\n-1\relax} {
      \draw[very thick, cap=round] (\x,0) -- (\x+1,0);
      \draw[very thick, cap=round] (\x,1) -- (\x+1,1);
    }
  \end{tikzpicture}%
}

% 2) dot squares (n x n)
\newcommand{\dotgrid}[2][black]{%
  \begin{tikzpicture}[scale=0.65]
    \pgfmathsetmacro\n{#2}
    \color{#1}
    \foreach \i in {0,...,\the\numexpr\n-1\relax} {
      \foreach \j in {0,...,\the\numexpr\n-1\relax} {
        \fill (0.4*\i,0.4*\j) circle (0.07);
      }
    }
  \end{tikzpicture}%
}

% 3) matchstick triangles (chain)
\newcommand{\matchtriangles}[2][black]{%
  \begin{tikzpicture}[scale=0.65]
    \pgfmathsetmacro\n{#2}
    \pgfmathsetmacro\h{0.866}
    \color{#1}
    \draw[very thick] (0,0) -- (\n,0);
    \draw[very thick] (0,0)
      \foreach \k in {1,...,\n} {
        -- ({\k - 0.5}, {\h}) -- (\k,0)
      };
  \end{tikzpicture}%
}

% 4) triangular numbers (dots)
\newcommand{\tridots}[2][black]{%
  \begin{tikzpicture}[scale=0.7]
    \pgfmathsetmacro\n{#2}
    \color{#1}
    \foreach \row in {1,...,\n} {
      \pgfmathsetmacro\k{int(\n - \row + 1)}
      \pgfmathsetmacro\xstart{ -0.45*(\k-1)/2 }
      \foreach \col in {0,...,\the\numexpr\k-1\relax} {
        \pgfmathsetmacro\x{\xstart + 0.45*\col}
        \pgfmathsetmacro\y{0.4*(\row-1)}
        \fill (\x,\y) circle (0.07);
      }
    }
  \end{tikzpicture}%
}

% 5) rectangle dots (n x (n+1))
\newcommand{\rectdots}[2][black]{%
  \begin{tikzpicture}[scale=0.65]
    \pgfmathsetmacro\n{#2}
    \color{#1}
    \foreach \i in {0,...,\the\numexpr\n-1\relax} {
      \foreach \j in {0,...,\n} {
        \fill (0.4*\i, 0.4*\j) circle (0.07);
      }
    }
  \end{tikzpicture}%
}

% 6) cross dots
\newcommand{\crossdots}[2][black]{%
  \begin{tikzpicture}[scale=0.65]
    \pgfmathsetmacro\n{#2}
    \color{#1}
    \fill (0,0) circle (0.07);
    \ifnum\n>1
      \foreach \d in {1,...,\the\numexpr\n-1\relax} {
        \fill (0, 0.4*\d) circle (0.07);
        \fill (0, -0.4*\d) circle (0.07);
        \fill (0.4*\d, 0) circle (0.07);
        \fill (-0.4*\d, 0) circle (0.07);
      }
    \fi
  \end{tikzpicture}%
}

\begin{document}

\frame{\titlepage}

% =============================================
% CONTENTS
% =============================================
\begin{frame}{Contents}
\small
\begin{columns}[T]
  \begin{column}{0.5\textwidth}
    \textbf{Questions and short answers}
    \begin{itemize}
      \item \hyperlink{q-st1}{\beamergotobutton{Starter 1}}
      \item \hyperlink{q-st2}{\beamergotobutton{Starter 2}}
      \item \hyperlink{q-st3}{\beamergotobutton{Starter 3}}
      \item \hyperlink{q-st4}{\beamergotobutton{Starter 4}}
      \item \hyperlink{q-st5}{\beamergotobutton{Starter 5}}
    \end{itemize}
  \end{column}
  \begin{column}{0.5\textwidth}
    \textbf{Worked solutions}
    \begin{itemize}
      \item \hyperlink{work-st1}{\beamergotobutton{Starter 1}}
      \item \hyperlink{work-st2}{\beamergotobutton{Starter 2}}
      \item \hyperlink{work-st3}{\beamergotobutton{Starter 3}}
      \item \hyperlink{work-st4}{\beamergotobutton{Starter 4}}
      \item \hyperlink{work-st5}{\beamergotobutton{Starter 5}}
    \end{itemize}
  \end{column}
\end{columns}
\end{frame}

% =============================================
% STARTER 1 – Easy
% =============================================
\begin{frame}{Starter 1}
  \hypertarget{q-st1}{}
  \begin{minipage}[t]{0.60\textwidth}
    \centering
    \textbf{Matchsticks}\\[0.7em]
    \matchsquares[red]{1}\;\hspace{0.5cm}
    \matchsquares[blue]{2}\;\hspace{0.5cm}
    \matchsquares[green!70!black]{3}\;\hspace{0.5cm}
    \textbf{?}\\[0.8em]
    \small
    1. Draw the next shape.\\[0.2em]
    2. How many matchsticks in the 4th?
  \end{minipage}\hfill
  \begin{minipage}[t]{0.34\textwidth}
    \centering
    \textbf{Dots}\\[0.7em]
    \dotgrid[red]{1}\;\hspace{0.5cm}
    \dotgrid[blue]{2}\;\hspace{0.5cm}
    \dotgrid[green!70!black]{3}\;\hspace{0.5cm}
    \textbf{?}\\[0.8em]
    \small
    3. Draw the next pattern.\\[0.2em]
    4. How many dots in the 4th?
  \end{minipage}

  \vspace{0.8cm}                     % increased vertical space
  \centering
  \textbf{5. Complete The Sequence}\\[0.3em]
  \small
  13,\;23,\;33,\;
  \ablank{43},\;
  \ablank{53},\;
  \ablank{63}
  \ashow{\ (add 10)}

  \vspace{1.2cm}                     % space before pattern answers
  \pause
  \begin{minipage}[t]{0.60\textwidth}
    \centering
    \textcolor{red}{\textbf{Answers}}\\[0.2em]
    \textcolor{red}{1. \matchsquares[orange]{4}}\\[0.2em]
    \textcolor{red}{2. 13 sticks}
  \end{minipage}\hfill
  \begin{minipage}[t]{0.34\textwidth}
    \centering
    \textcolor{red}{\textbf{Answers}}\\[0.2em]
    \textcolor{red}{3. \dotgrid[orange]{4}}\\[0.2em]
    \textcolor{red}{4. 16 dots}
  \end{minipage}
\end{frame}

% Worked solutions for Starter 1
\begin{frame}{Worked Solution: Starter 1, Question 1}
  \hypertarget{work-st1}{}
  \textbf{Question:} Draw the next matchstick shape.

  \textbf{Worked Solution:} The pattern has 1 square, then 2 squares, then 3 squares in a row. The next shape has 4 squares in a row.
  \begin{center}
    \matchsquares[orange]{4}
  \end{center}
\end{frame}

\begin{frame}{Worked Solution: Starter 1, Question 2}
  \textbf{Question:} How many matchsticks are in the 4th matchstick shape?

  \textbf{Worked Solution:} The counts are \(4, 7, 10, \dots\). Each extra square adds 3 matchsticks. Therefore shape 4 has
  \[
    10+3=13 \qquad 3(4)+1=13.
  \]
  So the 4th shape has 13 matchsticks.
\end{frame}

\begin{frame}{Worked Solution: Starter 1, Question 3}
  \textbf{Question:} Draw the next dot pattern.

  \textbf{Worked Solution:} The dot grids are \(1\times1\), then \(2\times2\), then \(3\times3\). The next pattern is \(4\times4\).
  \begin{center}
    \dotgrid[orange]{4}
  \end{center}
\end{frame}

\begin{frame}{Worked Solution: Starter 1, Question 4}
  \textbf{Question:} How many dots are in the 4th dot pattern?

  \textbf{Worked Solution:} The numbers of dots are \(1, 4, 9, \dots\), which are square numbers. Shape 4 has
  \[
    4^2 = 16
  \]
  dots.
\end{frame}

\begin{frame}{Worked Solution: Starter 1, Question 5}
  \textbf{Question:} Complete the sequence \(13, 23, 33, \_, \_, \_\).

  \textbf{Worked Solution:} The rule is add 10.
  \[
    33+10=43,\qquad 43+10=53,\qquad 53+10=63.
  \]
  Completed sequence: \(13, 23, 33, 43, 53, 63\).
\end{frame}

% =============================================
% STARTER 2 – Linear patterns
% =============================================
\begin{frame}{Starter 2}
  \hypertarget{q-st2}{}
  \begin{minipage}[t]{0.60\textwidth}
    \centering
    \textbf{Matchsticks}\\[0.7em]
    \matchtriangles[red]{1}\;\hspace{0.5cm}
    \matchtriangles[blue]{2}\;\hspace{0.5cm}
    \matchtriangles[green!70!black]{3}\;\hspace{0.5cm}
    \textbf{?}\\[0.8em]
    \small
    1. Draw shape 4.\\[0.2em]
    2. How many matchsticks in shape 4?
  \end{minipage}\hfill
  \begin{minipage}[t]{0.34\textwidth}
    \centering
    \textbf{Dots}\\[0.7em]
    \tridots[red]{1}\;\hspace{0.5cm}
    \tridots[blue]{2}\;\hspace{0.5cm}
    \tridots[green!70!black]{3}\;\hspace{0.5cm}
    \textbf{?}\\[0.8em]
    \small
    3. Draw pattern 4.\\[0.2em]
    4. How many dots in pattern 4?
  \end{minipage}

  \vspace{0.8cm}
  \centering
  \textbf{5. Complete The Sequence}\\[0.3em]
  \small
  1,\;3,\;5,\;7,\;
  \ablank{9},\;
  \ablank{11},\;
  \ablank{13}
  \ashow{\ (add 2)}

  \vspace{1.2cm}
  \pause
  \begin{minipage}[t]{0.60\textwidth}
    \centering
    \textcolor{red}{\textbf{Answers}}\\[0.2em]
    \textcolor{red}{1. \matchtriangles[orange]{4}}\\[0.2em]
    \textcolor{red}{2. 9 sticks}
  \end{minipage}\hfill
  \begin{minipage}[t]{0.34\textwidth}
    \centering
    \textcolor{red}{\textbf{Answers}}\\[0.2em]
    \textcolor{red}{3. \tridots[orange]{4}}\\[0.2em]
    \textcolor{red}{4. 10 dots}
  \end{minipage}
\end{frame}

% Worked solutions for Starter 2
\begin{frame}{Worked Solution: Starter 2, Question 1}
  \hypertarget{work-st2}{}
  \textbf{Question:} Draw shape 4 in the matchstick triangle pattern.

  \textbf{Worked Solution:} Each new shape adds one triangle to the row. Shape 4 therefore has 4 triangles.
  \begin{center}
    \matchtriangles[orange]{4}
  \end{center}
\end{frame}

\begin{frame}{Worked Solution: Starter 2, Question 2}
  \textbf{Question:} How many matchsticks are in shape 4?

  \textbf{Worked Solution:} The stick counts are \(3, 5, 7, \dots\) because each new triangle adds 2 sticks. So shape 4 has
  \[
    7+2=9.
  \]
  Equivalently, the rule is \(2n+1\), giving \(2(4)+1=9\).
\end{frame}

\begin{frame}{Worked Solution: Starter 2, Question 3}
  \textbf{Question:} Draw pattern 4 for the triangular dot pattern.

  \textbf{Worked Solution:} The triangular dot patterns have row lengths \(1;\; 1,2;\; 1,2,3\). Pattern 4 has row lengths \(1,2,3,4\).
  \begin{center}
    \tridots[orange]{4}
  \end{center}
\end{frame}

\begin{frame}{Worked Solution: Starter 2, Question 4}
  \textbf{Question:} How many dots are in pattern 4?

  \textbf{Worked Solution:} The dot counts are \(1, 3, 6, \dots\). Pattern 4 has
  \[
    1+2+3+4=10
  \]
  dots. Also, \(\frac{4\times5}{2}=10\).
\end{frame}

\begin{frame}{Worked Solution: Starter 2, Question 5}
  \textbf{Question:} Complete the sequence \(1, 3, 5, 7, \_, \_, \_.

  \textbf{Worked Solution:} Add 2 each time:
  \[
    7+2=9,\qquad 9+2=11,\qquad 11+2=13.
  \]
  Completed sequence: \(1, 3, 5, 7, 9, 11, 13\).
\end{frame}

% =============================================
% STARTER 3 – Predicting further terms
% =============================================
\begin{frame}{Starter 3}
  \hypertarget{q-st3}{}
  \begin{minipage}[t]{0.36\textwidth}
    \centering
    \textbf{Dots}\\[0.6em]
    \rectdots[red]{1}\;\hspace{0.5cm}
    \rectdots[blue]{2}\;\hspace{0.5cm}
    \rectdots[green!70!black]{3}\;\hspace{0.5cm}
    \textbf{?}\\[0.8em]
    \small
    1. Draw pattern 4.\\[0.2em]
    2. How many dots in pattern 4?
  \end{minipage}\hfill
  \begin{minipage}[t]{0.60\textwidth}
    \centering
    \textbf{Matchsticks}\\[0.7em]
    \matchsquares[red]{1}\;\hspace{0.5cm}
    \matchsquares[blue]{2}\;\hspace{0.5cm}
    \matchsquares[green!70!black]{3}\;\hspace{0.5cm}
    \textbf{?}\\[0.8em]
    \small
    3. How many sticks in the 10th shape?\\[0.2em]
    4. Find a rule.
  \end{minipage}

  \vspace{0.8cm}
  \centering
  \textbf{5. Complete The Sequence}\\[0.3em]
  \small
  3,\;6,\;9,\;12,\;
  \ablank{15},\;
  \ablank{18},\;
  \ablank{21}
  \ashow{\ (add 3)}
  \hspace{1em} \footnotesize 6. 20th term? \ashow{$60$}

  \vspace{0.5cm}
  \pause
  \begin{minipage}[t]{0.36\textwidth}
    \centering
    \textcolor{red}{\textbf{Answers}}\\[0.2em]
    \textcolor{red}{1. \rectdots[orange]{4}}\\[0.2em]
    \textcolor{red}{2. 20 dots \footnotesize (pattern is $n\times(n+1)$ dots)}
  \end{minipage}\hfill
  \begin{minipage}[t]{0.60\textwidth}
    \centering
    \textcolor{red}{\textbf{Answers}}\\[0.2em]
    \textcolor{red}{3. 31 }\\[0.2em]
    \textcolor{red}{4. Rule: $3n+1$}
  \end{minipage}
\end{frame}

% Worked solutions for Starter 3
\begin{frame}{Worked Solution: Starter 3, Question 1}
  \hypertarget{work-st3}{}
  \textbf{Question:} Draw pattern 4 for the rectangular dot pattern.

  \textbf{Worked Solution:} The rectangles are \(1\times2\), \(2\times3\), \(3\times4\), so pattern 4 is \(4\times5\).
  \begin{center}
    \rectdots[orange]{4}
  \end{center}
\end{frame}

\begin{frame}{Worked Solution: Starter 3, Question 2}
  \textbf{Question:} How many dots are in pattern 4?

  \textbf{Worked Solution:} Pattern 4 is a \(4\times5\) rectangle, so it has
  \[
    4\times5=20
  \]
  dots. In general, pattern \(n\) has \(n(n+1)\) dots.
\end{frame}

\begin{frame}{Worked Solution: Starter 3, Question 3}
  \textbf{Question:} How many sticks are in the 10th square-matchstick shape?

  \textbf{Worked Solution:} The square-matchstick counts are \(4, 7, 10, \dots\). The rule is \(3n+1\): shape \(n\) has \(n\) top sticks, \(n\) bottom sticks and \(n+1\) vertical sticks.
  \[
    3(10)+1=31.
  \]
  The 10th shape has 31 sticks.
\end{frame}

\begin{frame}{Worked Solution: Starter 3, Question 4}
  \textbf{Question:} Find a rule for the square-matchstick pattern.

  \textbf{Worked Solution:} If \(n\) is the shape number, the number of sticks is
  \[
    3n+1.
  \]
  Check: \(n=1\) gives 4, \(n=2\) gives 7, \(n=3\) gives 10.
\end{frame}

\begin{frame}{Worked Solution: Starter 3, Question 5}
  \textbf{Question:} Complete the sequence \(3, 6, 9, 12, \_, \_, \_.

  \textbf{Worked Solution:} Add 3 each time:
  \[
    12+3=15,\qquad 15+3=18,\qquad 18+3=21.
  \]
  Completed sequence: \(3, 6, 9, 12, 15, 18, 21\).
\end{frame}

\begin{frame}{Worked Solution: Starter 3, Question 6}
  \textbf{Question:} What is the 20th term of the sequence \(3, 6, 9, 12, \dots\)?

  \textbf{Worked Solution:} The sequence is the multiples of 3, so the \(n\)th term is \(3n\). Therefore the 20th term is
  \[
    3\times20=60.
  \]
\end{frame}

% =============================================
% STARTER 4 – Nth term rules
% =============================================
\begin{frame}{Starter 4}
  \hypertarget{q-st4}{}
  \begin{minipage}[t]{0.60\textwidth}
    \centering
    \textbf{Matchsticks}\\[0.7em]
    \matchtriangles[red]{1}\;\hspace{0.5cm}
    \matchtriangles[blue]{2}\;\hspace{0.5cm}
    \matchtriangles[green!70!black]{3}\;\hspace{0.5cm}
    \textbf{?}\\[0.8em]
    \small
    1. Find a rule for the number of matchsticks.\\[0.2em]
    2. Use it to find the 20th term.
  \end{minipage}\hfill
  \begin{minipage}[t]{0.34\textwidth}
    \centering
    \textbf{Dots}\\[0.7em]
    \tridots[red]{1}\;\hspace{0.5cm}
    \tridots[blue]{2}\;\hspace{0.5cm}
    \tridots[green!70!black]{3}\;\hspace{0.5cm}
    \textbf{?}\\[0.8em]
    \small
    3. Find a rule for the number of dots.\\[0.2em]
    4. Use it to find the 10th term.
  \end{minipage}

  \vspace{0.8cm}
  \centering
  \textbf{5. Complete The Sequence}\\[0.3em]
  \small
  5,\;9,\;13,\;17,\;
  \ablank{21},\;
  \ablank{25},\;
  \ablank{29}
  \ashow{\ (add 4)}
  \hspace{1em} \footnotesize 6. $n$th term rule? 100th term? \\ \hspace{12em} \ashow{$4n+1$, \quad \qquad \quad $401$}

  \vspace{1.2cm}
  \pause
  \begin{minipage}[t]{0.60\textwidth}
    \centering
    \textcolor{red}{\textbf{Answers}}\\[0.2em]
    \textcolor{red}{1. Rule: $2n+1$}\\[0.2em]
    \textcolor{red}{2. 20th term: $41$}
  \end{minipage}\hfill
  \begin{minipage}[t]{0.34\textwidth}
    \centering
    \textcolor{red}{\textbf{Answers}}\\[0.2em]
    \textcolor{red}{3. Rule: $\frac{n(n+1)}{2}$}\\[0.2em]
    \textcolor{red}{4. 10th term: $55$}
  \end{minipage}
\end{frame}

% Worked solutions for Starter 4
\begin{frame}{Worked Solution: Starter 4, Question 1}
  \hypertarget{work-st4}{}
  \textbf{Question:} Find a rule for the number of matchsticks in the triangle-stick pattern.

  \textbf{Worked Solution:} The counts are \(3, 5, 7, \dots\). Each term increases by 2, so the rule is
  \[
    2n+1.
  \]
  Check: \(n=1\) gives 3, \(n=2\) gives 5, \(n=3\) gives 7.
\end{frame}

\begin{frame}{Worked Solution: Starter 4, Question 2}
  \textbf{Question:} Use the rule to find the 20th term.

  \textbf{Worked Solution:} Using \(2n+1\) with \(n=20\),
  \[
    2(20)+1=41.
  \]
  The 20th term is 41.
\end{frame}

\begin{frame}{Worked Solution: Starter 4, Question 3}
  \textbf{Question:} Find a rule for the number of dots in the triangular dot pattern.

  \textbf{Worked Solution:} The dot counts are triangular numbers \(1, 3, 6, \dots\). Pattern \(n\) has
  \[
    1+2+\cdots+n=\frac{n(n+1)}{2}
  \]
  dots.
\end{frame}

\begin{frame}{Worked Solution: Starter 4, Question 4}
  \textbf{Question:} Use the rule to find the 10th term.

  \textbf{Worked Solution:} Using \(\frac{n(n+1)}{2}\) with \(n=10\),
  \[
    \frac{10(10+1)}{2}=\frac{110}{2}=55.
  \]
  The 10th term is 55.
\end{frame}

\begin{frame}{Worked Solution: Starter 4, Question 5}
  \textbf{Question:} Complete the sequence \(5, 9, 13, 17, \_, \_, \_.

  \textbf{Worked Solution:} Add 4 each time:
  \[
    17+4=21,\qquad 21+4=25,\qquad 25+4=29.
  \]
  Completed sequence: \(5, 9, 13, 17, 21, 25, 29\).
\end{frame}

\begin{frame}{Worked Solution: Starter 4, Question 6}
  \textbf{Question:} For the sequence \(5, 9, 13, 17, \dots\), what is the \(n\)th term rule and the 100th term?

  \textbf{Worked Solution:} The first term is 5 and the common difference is 4, so
  \[
    T_n = 5 + 4(n-1) = 4n+1.
  \]
  The 100th term is
  \[
    4(100)+1=401.
  \]
\end{frame}

% =============================================
% STARTER 5 – Challenge
% =============================================
\begin{frame}{Starter 5}
  \hypertarget{q-st5}{}
  \begin{minipage}[t]{0.50\textwidth}
    \centering
    \textbf{Dots}\\[0.7em]
    \crossdots[red]{1}\;\hspace{0.5cm}
    \crossdots[blue]{2}\;\hspace{0.5cm}
    \crossdots[green!70!black]{3}\;\hspace{0.5cm}
    \textbf{?}\\[0.8em]
    \small
    1. Draw shape 4.\\[0.2em]
    2. How many dots in shape 5?\\[0.2em]
    3. Try to find a rule.
  \end{minipage}\hfill
  \begin{minipage}[t]{0.50\textwidth}
    \centering
    \textbf{Special Sequence}\\[0.7em]
    1,\;1,\;2,\;3,\;5,\;8,\;
    \ablank{13},\;
    \ablank{21},\;
    \ablank{34}\\[0.8em]
    \small
    4. Next three terms?\\[0.2em]
    5. What is the rule?
  \end{minipage}

  \vspace{0.8cm}
  \centering
  \textbf{6. Complete The Sequence}\\[0.3em]
  \small
  1,\;4,\;9,\;16,\;
  \ablank{25},\;
  \ablank{36},\;
  \ablank{49}
  \ashow{\ (square numbers $n^2$)}
  \hspace{1em} \footnotesize 7. 20th term? \ashow{$400$}

  \vspace{0.5cm}
  \pause
  \begin{minipage}[t]{0.50\textwidth}
    \begin{multicols}{2}
    \textcolor{red}{\textbf{Answers}}\\[0.2em]
    \textcolor{red}{1. \crossdots[orange]{4}}\\[0.2em]
    \textcolor{red}{2. Shape 5 has 17 dots}\\[0.2em]
    \textcolor{red}{3. Rule: $4n-3$}
    \end{multicols}
  \end{minipage}\hfill
  \begin{minipage}[t]{0.50\textwidth}
    \centering
    \vspace{-6.5em}
    \textcolor{red}{\textbf{Answers}}\\[0.2em]
    \textcolor{red}{4. 13,\;21,\;34}\\[0.2em]
    \textcolor{red}{5. Fibonacci: add the two previous terms}
  \end{minipage}
\end{frame}

% Worked solutions for Starter 5
\begin{frame}{Worked Solution: Starter 5, Question 1}
  \hypertarget{work-st5}{}
  \textbf{Question:} Draw shape 4 for the cross-dot pattern.

  \textbf{Worked Solution:} Each new shape adds one dot to each of the four arms. Shape 4 has a centre dot and 3 dots on each arm.
  \begin{center}
    \crossdots[orange]{4}
  \end{center}
\end{frame}

\begin{frame}{Worked Solution: Starter 5, Question 2}
  \textbf{Question:} How many dots are in shape 5?

  \textbf{Worked Solution:} Shape 5 has 1 centre dot plus 4 dots on each of 4 arms:
  \[
    1+4\times4=17.
  \]
  Equivalently, using the rule \(4n-3\), \(4(5)-3=17\).
\end{frame}

\begin{frame}{Worked Solution: Starter 5, Question 3}
  \textbf{Question:} Try to find a rule for the cross-dot pattern.

  \textbf{Worked Solution:} Shape \(n\) has a centre dot and four arms of length \(n-1\). Therefore
  \[
    1+4(n-1)=4n-3.
  \]
  So the rule is \(4n-3\).
\end{frame}

\begin{frame}{Worked Solution: Starter 5, Question 4}
  \textbf{Question:} Next three terms in \(1, 1, 2, 3, 5, 8, \_, \_, \_\).

  \textbf{Worked Solution:} Add the two previous terms:
  \[
    5+8=13,\qquad 8+13=21,\qquad 13+21=34.
  \]
  Next three terms: \(13, 21, 34\).
\end{frame}

\begin{frame}{Worked Solution: Starter 5, Question 5}
  \textbf{Question:} What is the rule for the special sequence?

  \textbf{Worked Solution:} Start with 1 and 1. After that, each term is the sum of the two previous terms. This is the Fibonacci sequence.
\end{frame}

\begin{frame}{Worked Solution: Starter 5, Question 6}
  \textbf{Question:} Complete the sequence \(1, 4, 9, 16, \_, \_, \_.

  \textbf{Worked Solution:} These are square numbers:
  \[
    5^2=25,\qquad 6^2=36,\qquad 7^2=49.
  \]
  Completed sequence: \(1, 4, 9, 16, 25, 36, 49\).
\end{frame}

\begin{frame}{Worked Solution: Starter 5, Question 7}
  \textbf{Question:} What is the 20th term of \(1, 4, 9, 16, \dots\)?

  \textbf{Worked Solution:} The \(n\)th term is \(n^2\). Therefore the 20th term is
  \[
    20^2=400.
  \]
\end{frame}

\end{document}